\clearpage
\part{第一题：基础实现}
\section{程序基本功与做题技巧}
\subsection{程序基本功}
程序基本功包括C语言的格式化输入输出、整型与浮点运算、函数的编写与调用、数组和结构体的使用。
\subsubsection{格式化输入输出}
scanf把输入存入变量，普通变量前写地址符\&；printf输出变量的值。格式符必须与类型匹配。
\begin{center}
\begin{tabular}{llll}
\toprule 类型 & 示例声明 & scanf & printf\\\midrule
int & \texttt{int n;} & \texttt{\%d} & \texttt{\%d}\\
long long & \texttt{long long s;} & \texttt{\%lld} & \texttt{\%lld}\\
double & \texttt{double x;} & \texttt{\%lf} & \texttt{\%.6f}（6位小数）\\
char & \texttt{char c;} & \texttt{\%c} & \texttt{\%c}\\
字符数组 & \texttt{char s[101];} & \texttt{\%100s} & \texttt{\%s}\\
\bottomrule
\end{tabular}
\end{center}
\begin{lstlisting}
int n; long long sum; double x; char c;
scanf("%d%lld%lf %c", &n, &sum, &x, &c);
printf("%d %lld %.6f %c\n", n, sum, x, c);
char word[101];
scanf("%100s", word);
\end{lstlisting}
\texttt{\%c}会读取空白，前面的空格表示先跳过空白。\texttt{\%s}读到空白就停止；上例最多读100个字符，另留一格存结束符。字符数组名表示首元素地址，不再加\&。
\subsubsection{整型与浮点运算}
整数除法向0截断，\texttt{5/2}得到2，\texttt{-5/2}得到$-2$；\texttt{\%}取整数余数。要保留小数，至少把一个操作数转为浮点数。
\begin{lstlisting}
int a = 1000000000, b = 3;
long long product = 1LL * a * b;
double mean = (double)a / b;
int ceiling = a / b + (a % b != 0);
\end{lstlisting}
乘法先按操作数类型计算，再赋值；\texttt{1LL}使乘法按long long进行。最后一行适用于$a\ge0,b>0$，计算向上取整的商。double能表示小数，但有舍入误差；输出精度由printf格式控制。
\subsubsection{函数}
函数用参数接收数据，用返回值给出结果。参数按值传递；要修改调用方的变量，可以传入它的地址。数组参数需要同时给出长度。
\begin{lstlisting}
long long sum_array(const int a[], int n) {
    long long s = 0;
    for (int i = 0; i < n; i++) s += a[i];
    return s;
}
void swap_int(int *x, int *y) {
    int t = *x; *x = *y; *y = t;
}
// Calls: sum_array(a, n); swap_int(&a[0], &a[1]);
\end{lstlisting}
const表示函数不通过a修改数组。\texttt{*x}访问地址x处的变量，所以交换结果会保留在原数组中。
\subsubsection{数组}
数组把同类型元素连续存放，\texttt{a[i]}访问第i格。长度为n时，合法下标是0到$n-1$；二维数组\texttt{grid[i][j]}分别用行、列下标访问。
\begin{lstlisting}
const int N = 1005;
int a[N], grid[N][N]; // Global: initially zero.
int count[10] = {0};  // Local: explicitly initialized.
for (int i = 0; i < n; i++) scanf("%d", &a[i]);
\end{lstlisting}
全局数组自动清零，未初始化的局部数组不能直接读取。频率数组\texttt{count[x]}按数值x计数，要求x处于数组的合法下标范围。
\subsubsection{结构体}
结构体把同一对象的多个属性放在一起，通过点号访问成员；多个对象可以存入结构体数组。
\begin{lstlisting}
struct Point { int x, y; char type; };
struct Point p[1005];
scanf("%d%d %c", &p[0].x, &p[0].y, &p[0].type);
int distance2 = p[0].x * p[0].x + p[0].y * p[0].y;
\end{lstlisting}
最后一行计算点到原点的距离平方，假定坐标绝对值不超过1000，此时int足够。
\subsection{做题技巧}
\subsubsection{数组容量与运算范围}
数组容量由最大访问下标决定。若$n\le1000$且使用1下标，至少开1001格；若还访问$n+1$，至少开1002格。字符串长度最多为L时，字符数组至少开$L+1$格。

常用评测环境中int为32位，范围约为$\pm2.14\times10^9$；long long为64位，范围约为$\pm9.22\times10^{18}$。把数据范围代入每一步运算：$2\times10^5$项、每项绝对值至多$10^9$，和的绝对值至多$2\times10^{14}$，需要long long；$10^{18}\times100$则连long long也放不下。
\subsubsection{时间复杂度与估算}
时间复杂度表示操作次数随输入规模增长的数量级。先数循环执行次数，再忽略常数倍和低阶项；$O(n)$表示线性增长，$O(n^2)$表示平方增长。

\Needspace{8\baselineskip}\textbf{一次遍历：$O(n)$。}
\begin{lstlisting}
long long sum = 0;
for (int i = 0; i < n; i++) sum += a[i];
\end{lstlisting}
执行n次加法。n扩大10倍，操作次数也扩大10倍。

\Needspace{8\baselineskip}\textbf{两次依次执行的遍历：$O(n)$。}
\begin{lstlisting}
long long sum = 0; int positive = 0;
for (int i = 0; i < n; i++) sum += a[i];
for (int i = 0; i < n; i++) if (a[i] > 0) positive++;
\end{lstlisting}
两个循环各执行n次，合计约2n次，忽略常数2后仍是$O(n)$。

\Needspace{9\baselineskip}\textbf{枚举所有数对：$O(n^2)$。}
\begin{lstlisting}
long long pairs = 0;
for (int i = 0; i < n; i++)
    for (int j = i + 1; j < n; j++)
        if (a[i] == a[j]) pairs++;
\end{lstlisting}
比较次数是$(n-1)+(n-2)+\cdots+1=n(n-1)/2$。n扩大10倍，次数约扩大100倍。

\Needspace{8\baselineskip}\textbf{遍历n行m列：$O(nm)$。}
\begin{lstlisting}
long long sum = 0;
for (int i = 0; i < n; i++)
    for (int j = 0; j < m; j++) sum += grid[i][j];
\end{lstlisting}
共nm次加法；若m始终是固定常数，便可记作$O(n)$。

\Needspace{8\baselineskip}\textbf{每次减半：$O(\log n)$。}
\begin{lstlisting}
int steps = 0;
for (int length = n; length > 1; length /= 2) steps++;
\end{lstlisting}
n=16时依次处理16、8、4、2，执行4次。通常约为$\log_2 n$次，n翻倍只多一次。

将最大规模代入估算：$n=10^5$时，遍历约$10^5$次，数对枚举约$5\times10^9$次，后者通常无法在几秒内完成。实际耗时还受操作类型、输入输出和机器速度影响。
\subsubsection{空间复杂度与估算}
空间复杂度表示存储量的增长数量级：n项数组为$O(n)$，$n\times m$矩阵为$O(nm)$，固定数量的变量为$O(1)$。深度n的递归还需要$O(n)$调用栈空间。

实际内存约为元素个数乘以每个元素的字节数，再加上其他数组、容器和调用栈。$1000\times1000$的int矩阵约3.8MiB，$10^5$项long long数组约0.76MiB。大数组通常放在全局，避免局部栈容量不足；结构体大小可用sizeof查询。
